% Exact quarantined release candidate for P06.2.
% This file is not a live-book input before independent and production
% verification.  No novelty or priority claim is made.

\begin{proof}[Proof of \cref{prop:p062-rational-centred-traces}]
Fix \(a\in D^-_{\mathbf R}(K)\) and a polyradius
\(\mathbf r\) with \(0<r_i<R_i\).  Choose numbers
\begin{equation}\label{eq:p062-rho-above-r}
                         r_i<\rho_i<R_i
\end{equation}
for every \(i\), and put
\[
 \Omega_{\boldsymbol\rho}
   =\prod_{i=1}^n(-\log\rho_i,\infty),
 \qquad
 A_{a,\boldsymbol\rho}
   =\{0<|T_i-a_i|<\rho_i\text{ for every }i\}.
\]
The ultrametric inequality places \(A_{a,\boldsymbol\rho}\) inside
\(D^-_{\mathbf R}\).  Translation by \(a\) identifies it with the split
torus domain
\(\operatorname{trop}^{-1}(\Omega_{\boldsymbol\rho})\), whose
tropicalization and weight-one Gauss section are
\[
 F_a=(-\log|T_1-a_1|,\ldots,-\log|T_n-a_n|),
 \qquad
 \iota_a(u)=\eta_{a,(\mathrm e^{-u_1},\ldots,\mathrm e^{-u_n})}.
\]

Let \(\beta\) be a compactly supported smooth Lagerberg form of
bidegree \((n-1,n-1)\) on \(\Omega_{\boldsymbol\rho}\).  Tropicalization
of a split torus is proper.  Thus \(F_a^*\beta\) is a compactly supported
smooth test form on \(A_{a,\boldsymbol\rho}\).  Restricting the current
equation to this ordinary open analytic subspace, or equivalently
extending this test form by zero, gives
\begin{equation}\label{eq:p062-arbitrary-field-test-identity}
 \begin{aligned}
 0
 &=\bigl\langle\mathrm d'\mathrm d''[h],F_a^*\beta\bigr\rangle\\
 &=\int_{A_{a,\boldsymbol\rho}}
             h\,F_a^*(\mathrm d'\mathrm d''\beta).
 \end{aligned}
\end{equation}

Set \(\alpha=\mathrm d'\mathrm d''\beta\) and let
\(\mu_\alpha=(F_a^*\alpha)_{\mathrm{Rad}}\) be its signed Radon
measure.  Choose a compact \((\mathbf Z,-\log|K^\times|)\)-polytope
\(\Delta\Subset\Omega_{\boldsymbol\rho}\) containing
\(\operatorname{supp}(\alpha)\).  Such a polytope exists even when the
valuation is discrete: the divisible hull of the nonzero additive group
\(-\log|K^\times|\) is dense in \(\mathbf R\), and rational integral
inequalities therefore provide a compact box strictly inside
\(\Omega_{\boldsymbol\rho}\) around the given support.  The top-form
support theorem on a torus
\cite[Proposition~2.3.6]{BGJKK25} and the weight-one calibrated
integration formula
\cite[Proposition~11.1 and Example~11.7]{GJR26} give
\begin{equation}\label{eq:p062-arbitrary-field-support-push}
 \operatorname{supp}(\mu_\alpha)
       \subseteq\iota_a(\operatorname{supp}\alpha),
 \qquad
 (F_a)_*\mu_\alpha=\alpha_{\mathrm{Rad}}.
\end{equation}
These statements do not require \(K\) to be algebraically closed.  The
calibrated formula first gives the second equality against smooth
compactly supported multipliers.  Such multipliers are uniformly dense
on the common compact support, so uniqueness of signed Radon measures
gives the displayed equality.  Since \(F_a\circ\iota_a\) is the identity,
the support and pushforward identities imply
\begin{equation}\label{eq:p062-arbitrary-field-measure-section}
                      \mu_\alpha
                         =(\iota_a)_*\alpha_{\mathrm{Rad}}.
\end{equation}
The multiplier \(h\) is continuous, so
\eqref{eq:p062-arbitrary-field-test-identity} and
\eqref{eq:p062-arbitrary-field-measure-section} yield
\begin{equation}\label{eq:p062-arbitrary-field-trace-hessian}
  \int_{\Omega_{\boldsymbol\rho}}
       h(\iota_a(u))\,\mathrm d'\mathrm d''\beta(u)=0
\end{equation}
for every such \(\beta\).  Thus the continuous function
\[
                    \tau_a(u)=h(\iota_a(u))
\]
has zero distributional Lagerberg Hessian.  In coordinates, every
distribution \(\partial^2\tau_a/\partial u_i\partial u_j\) vanishes.
Convolution on relatively compact Euclidean balls shows that every local
regularization is affine; compatibility on overlapping balls and passage
to the continuous limit give
\begin{equation}\label{eq:p062-arbitrary-field-affine-trace}
                    \tau_a(u)=c+\sum_{i=1}^n\lambda_i u_i.
\end{equation}

Fix \(i\) and all \(u_j\) with \(j\ne i\).  As \(u_i\to+\infty\), the
Gauss seminorms converge to the Gauss seminorm on the coordinate section
\(T_i=a_i\).  This follows directly from
\[
 \left|\sum_{\nu\in\mathbf N^n}c_\nu(T-a)^\nu\right|_{
      \eta_{a,(\mathrm e^{-u_1},\ldots,\mathrm e^{-u_n})}}
       =\max_\nu |c_\nu|\mathrm e^{-\langle\nu,u\rangle},
\]
first for polynomials and then, by uniform approximation, on a fixed
closed polydisc.  The limiting point lies in \(D^-_{\mathbf R}\), and
continuity of \(h\) gives a finite limit.  Hence \(\lambda_i=0\) in
\eqref{eq:p062-arbitrary-field-affine-trace}.  Repeating this for every
\(i\) shows that \(\tau_a\) is constant.  Letting all coordinates tend
to \(+\infty\) gives its value \(h(a)\).  Because of
\eqref{eq:p062-rho-above-r}, the point corresponding to the originally
chosen \(\mathbf r\) belongs to this trace.  This proves
\eqref{eq:p062-rational-centred-radial-identity} for every allowed
\(\mathbf r\).

If \(a,b\in D^-_{\mathbf R}(K)\), choose
\[
                  |a_i-b_i|\leq r_i<\rho_i<R_i.
\]
Translation invariance of the Gauss norm gives
\(\eta_{a,\mathbf r}=\eta_{b,\mathbf r}\).  The radial identity at the
two centres gives \(h(a)=h(b)\).  Hence the value on every \(K\)-centred
Gauss skeleton is the same, and continuity gives constancy on the closure
of their union.

Now let \(h\) be defined on \(D^+_{\mathbf R}\) and suppose that the
current equation holds on its absolute analytic interior.  For a
\(K\)-point \(a\) and radii \(r_i<\rho_i<R_i\), the set
\(A_{a,\boldsymbol\rho}\) is an ambient open analytic subspace contained
in \(D^+_{\mathbf R}\).  It is therefore contained in the topological
interior of that analytic domain.  Berkovich's relative-interior theorem
and composition formula
\cite[Proposition~1.5.5(i)--(ii)]{Berk93}, together with the
boundarylessness of analytic affine space, place it in the absolute
analytic interior.  The preceding test-form argument applies and proves
the radial identity when every \(r_i<R_i\).  Letting any selected radii
increase to \(R_i\), and using continuity of the Gauss section and of
\(h\), proves it for \(r_i\leq R_i\).  For two \(K\)-points \(a,b\), take
\(0<r_i\leq R_i\) with
\(|a_i-b_i|\leq r_i\), using \(r_i=R_i\) when necessary.  The same common
Gauss point argument finishes the closed case.
\end{proof}

\begin{proof}[Proof of \cref{prop:p062-explicit-base-change}]
Let \(p_L:D^-_{\mathbf R,L}\to D^-_{\mathbf R}\) be the scalar-extension
projection.  Apply \cref{prop:p062-rational-centred-traces} over \(L\) to
\(h_L=h\circ p_L\).  It follows that \(h_L\) is constant on the
\(L\)-rational points.  Rigid points are dense in every strictly affinoid
space \cite[Proposition~2.1.15]{Berk12}; because \(L\) is algebraically
closed, the rigid points of a strict \(L\)-polydisc are precisely its
\(L\)-rational points.  The open polydisc is exhausted pointwise by strict
closed subpolydiscs: the nontrivial divisible group
\(|L^\times|\) is dense in \(\mathbf R_{>0}\), so the subpolydisc radii
may be chosen strictly between the coordinate seminorms of a given point
and the \(R_i\).  Indeed, after applying \(-\log\), this value group is
a nonzero divisible additive subgroup and contains
\(\gamma\mathbf Q\) for some \(\gamma\ne0\).  Therefore the
\(L\)-rational points are dense in
\(D^-_{\mathbf R,L}\), and continuity makes \(h_L\) constant everywhere.
The canonical projection \(p_L\) is surjective
\cite[\S11.1.9]{CLD25}, so \(h\) is constant.
This proves the first item.  The strict closed polydisc is already
affinoid; the closed part of \cref{prop:p062-rational-centred-traces},
the same rigid-point density, and surjectivity of scalar extension prove
the second item.

Assume now that both \(h\) and \(-h\) are CLD-psh in the approximation
sense.  The scalar-extension projection is a morphism of the associated
affined G-tropical spaces: pullback carries every local affine tropical
function to a local affine tropical function
\cite[\S\S12.4.5--12.4.6]{CLD25}.  Functoriality of psh
functions \cite[Proposition~8.1.5]{CLD25} therefore makes \(h_L\) and
\(-h_L\) psh.  The open polydisc is separated, pure \(n\)-dimensional,
and boundaryless, and the two continuous functions are locally
integrable.  Thus their Hessian currents are positive by
\cite[Proposition~8.3.1(a)]{CLD25}.  Since
\[
             \mathrm d'\mathrm d''[-h_L]
                 =-\mathrm d'\mathrm d''[h_L],
\]
both a current and its negative are positive.  Evaluation on positive
test forms shows that both currents vanish.  The first item now applies.
\end{proof}

% Research candidate for P06.2.  Not authorized for live-book use before
% independent exact-byte audit and production Rethlas verification.

\begin{proof}[Proof of \cref{prop:p062-finite-flat-continuous-trace}]
The assertions are local on \(X\), so fix an affinoid domain \(U\subset X\).
Its inverse image \(Y_U\) is affinoid.  By the smooth
Stone--Weierstrass theorem of Chambert--Loir--Ducros
\cite[Theorem~15.2.6]{CLD25}, there are ordinary CLD-smooth functions
\(h_j\) on \(Y_U\) which converge uniformly to \(h|_{Y_U}\).
For every \(j\), put
\[
 H_j(x):=\int_{Y_x}h_j|_{Y_x},\qquad x\in U.
\]
By \cite[Theorem~16.3.3(d)]{CLD25}, fibre integration in relative
dimension zero makes \(H_j\) \(G\)-smooth, hence continuous, on \(U\).
Every fibre has
calibrated mass \(d\).  Consequently
\begin{equation}\label{eq:p062-trace-uniform-bound}
 \lVert H_j-\operatorname{Tr}_p(h)\rVert_{C^0(U)}
 \leq d\,\lVert h_j-h\rVert_{C^0(Y_U)}.
\end{equation}
Thus the trace in \eqref{eq:p062-finite-trace-definition} is continuous.

For a smooth \(h_j\), the projection formula for fibre integration
\cite[Corollary~16.5.5]{CLD25} gives
\begin{equation}\label{eq:p062-smooth-trace-current}
                         p_*[h_j]=[H_j]
\end{equation}
on \(U\).  Uniform convergence against every compactly supported smooth
top-degree test form lets \(j\to\infty\) in
\eqref{eq:p062-smooth-trace-current}.  Together with
\eqref{eq:p062-trace-uniform-bound}, this proves
\eqref{eq:p062-finite-trace-current} on \(U\), and hence on \(X\).

Finally, direct image by a compact morphism commutes with
\(\mathrm d'\) and \(\mathrm d''\)
\cite[Section~7.7.8]{CLD25}.  Restricting
\eqref{eq:p062-finite-trace-current} to \(V\) therefore gives
\[
 \mathrm d'\mathrm d''[\operatorname{Tr}_p(h)]
   =p_*\bigl(\mathrm d'\mathrm d''[h]\bigr)=0.
\]
This proves \eqref{eq:p062-finite-trace-hessian}.
\end{proof}

\begin{proof}[Proof of \cref{prop:p062-arbitrary-field-disc}]
We treat the open and closed discs simultaneously until the final density
step.  Let \(y\) be a rigid point of the disc under consideration, let
\(E=\mathscr H(y)\), and let \(P\in K[T]\) be the monic irreducible
polynomial of \(T(y)\).  Write
\[
 P(T)=T^d+\sum_{j=0}^{d-1}c_jT^j,
 \qquad d=[E:K],
 \qquad s=|T(y)|.
\]
The absolute value of \(K\) has a unique extension to every finite
extension.  Hence every algebraic conjugate of \(T(y)\), counted with its
inseparable multiplicity, has absolute value \(s\).  The elementary
symmetric functions of these conjugates give
\begin{equation}\label{eq:p062-minimal-polynomial-coefficient-bound}
                         |c_j|\leq s^{d-j}
                         \qquad(0\leq j<d).
\end{equation}

In the open case one has \(s<R\), and the monic polynomial \(P\) induces
a finite flat morphism of degree \(d\)
\begin{equation}\label{eq:p062-open-polynomial-map}
 p\colon D_R^-\longrightarrow D_{R^d}^- ,
 \qquad Z\longmapsto P(T).
\end{equation}
Indeed, \eqref{eq:p062-minimal-polynomial-coefficient-bound} shows that
\(|P(x)|<R^d\) if and only if \(|T(x)|<R\).  The morphism is the indicated
base change of the finite free polynomial morphism
\(\mathbf A_K^1\to\mathbf A_K^1\); freeness, with basis
\(1,T,\ldots,T^{d-1}\), also follows directly by division by the monic
polynomial \(P(T)-Z\).

In the closed case, \(s\leq R\), and the same argument gives a finite flat
degree-\(d\) morphism
\begin{equation}\label{eq:p062-closed-polynomial-map}
 p\colon D_R^+\longrightarrow D_{R^d}^+ .
\end{equation}
The fibre over the \(K\)-rational point \(0\) is
\(
 \mathcal M(K[T]/(P))=\{y\}
\)
and has calibrated degree \(d\).

Set \(H=\operatorname{Tr}_p(h)\).  By
\cref{prop:p062-finite-flat-continuous-trace}, \(H\) is continuous.  In the
open case it satisfies
\(
 \mathrm d'\mathrm d''[H]=0
\)
on \(D_{R^d}^-\).  In the closed case the same equation holds on the
absolute analytic interior of \(D_{R^d}^+\).  For the latter assertion,
the Gauss-fibre calculation below, applied with \(t=R\), shows that the
inverse image of the unique analytic boundary point \(\xi_{R^d}\) is the
unique analytic boundary point \(\eta_R\).  Thus the pullback of a test form
supported in the target interior is supported in the source interior, where
the current equation for \(h\) holds.

We next record the fibre at a suitable Gauss point.  In the open case choose
\begin{equation}\label{eq:p062-radius-above-rigid-point}
                              s<t<R,
\end{equation}
and put \(q=t^d\).  In the closed case make the same choice when \(s<R\),
and put \(t=R\), \(q=R^d\), when \(s=R\).  Let \(\eta_t\) be the Gauss
point of the source disc and \(\xi_q\) the Gauss point of the target disc.
Then
\begin{equation}\label{eq:p062-single-gauss-fibre}
                         p^{-1}(\xi_q)=\{\eta_t\},
 \qquad \deg_{\eta_t}(p)=d.
\end{equation}

For completeness, \eqref{eq:p062-single-gauss-fibre} follows directly from
Gauss valuations.  By
\eqref{eq:p062-minimal-polynomial-coefficient-bound}, the \(t\)-Gauss norm
of \(P\) is \(q=t^d\), and its initial form has degree \(d\).  If
\(U\) and \(V\) denote the homogeneous residues of \(T\) and \(Z\), the
initial equation of \(P(T)-Z\) is
\[
                      \operatorname{in}_t(P)(U)-V.
\]
The element \(\operatorname{in}_t(P)(U)\) is transcendental over the
graded residue field.  It follows first that composition with \(P\) carries
the \(t\)-Gauss valuation to the \(q\)-Gauss valuation, so
\(p(\eta_t)=\xi_q\).  Moreover, \(V\) is transcendental over the graded
residue field, and the displayed
polynomial is irreducible over its graded fraction field: equivalently,
the rational-function extension generated by \(U\) over that generated by
\(\operatorname{in}_t(P)(U)\) has degree \(d\).  This remains true in the
inseparable case.  The graded fundamental inequality therefore gives local
degree at least \(d\) at \(\eta_t\).  Since the finite flat morphism has total
degree \(d\), there is no other point in the fibre and the local degree is
exactly \(d\).  This proves \eqref{eq:p062-single-gauss-fibre} without
splitting \(P\) into geometric roots.

The zero fibre and \eqref{eq:p062-single-gauss-fibre} now give
\begin{equation}\label{eq:p062-trace-two-fibres}
                         H(0)=d\,h(y),
 \qquad                  H(\xi_q)=d\,h(\eta_t).
\end{equation}
Apply \cref{prop:p062-rational-centred-traces} to \(H\) on the target
disc.  Its radial identity, including the endpoint in the closed case,
and \eqref{eq:p062-trace-two-fibres} yield
\begin{equation}\label{eq:p062-rigid-to-gauss}
                              h(y)=h(\eta_t).
\end{equation}
Applying the same proposition to \(h\) on the source disc gives
\(h(\eta_t)=h(0)\).  Thus every rigid point \(y\) has the same value
\(h(0)\).

Rigid points are dense in every strict affinoid domain; see
\cite[Proposition~2.1.15]{Berk12}.  This applies directly to \(D_R^+\).
Applied locally on a strict-affinoid covering of \(D_R^-\), it also proves
density there.  Continuity of \(h\) proves constancy on the entire open or
closed disc.
\end{proof}

% Exact frozen candidate proof: one-dimensional endpoint variants only.

\begin{proof}
For \(s\in\{r,R\}\), let \(\eta_s\) denote the \(s\)-Gauss point.
The absolute interior of \(A\) is obtained by deleting precisely the Gauss
point corresponding to each included endpoint.  In particular it contains
the strict radial core
\[
                         A^-_{r,R}=\{r<|T|<R\}
\]
and every nonskeletal residue branch at an included endpoint.

Here are the precise topological inputs.  Over an algebraically closed
complete field, the decomposition of a generalized annulus off its skeleton
into open balls, together with the closure of each ball at its retraction
point, is \cite[Lemma~2.12]{BPR13Structure}; that paper assumes algebraic
closedness throughout.  Over an arbitrary complete field, the corresponding
components are virtual discs by
\cite[\S\S3.6.35.5--3.6.35.10, especially 3.6.35.8]
{Ducros24CurveStructure}, and the skeleton and retraction commute with
complete scalar extension.  Berkovich's relative-interior theorem and
composition formula \cite[Proposition~1.5.5(i)--(ii)]{Berk93}, applied in the
boundaryless analytic torus, identify the relative topological interior with
the absolute analytic interior.  A closed endpoint Gauss point is not in that
interior because the ambient torus skeleton continues through it, whereas
every nonskeletal endpoint branch is ambient open.  This proves the asserted
description of \(\operatorname{Int}(A)\).

Restriction of \eqref{eq:p062-annulus-endpoint-zero-current} to
\(A^-_{r,R}\), followed by the arbitrary-field strict-open annulus theorem,
gives unique \(a,b\in\mathbf R\) such that
\begin{equation}\label{eq:p062-annulus-endpoint-core-formula}
                         h=a\log|T|+b
                         \qquad\text{on }A^-_{r,R}.
\end{equation}
Continuity along the annular skeleton gives the same formula at every
included endpoint Gauss point.  It remains to prove it on the nonskeletal
endpoint branches.

Fix an included endpoint \(s\in\{r,R\}\), and let \(y\in A\) be a rigid
point with \(|T(y)|=s\).  Let
\[
                  P(T)=T^d+\sum_{i=0}^{d-1}c_iT^i
\]
be the monic irreducible polynomial of \(T(y)\) over \(K\).  Completeness of
\(K\) gives a unique extension of its absolute value to every finite
extension.  Hence all roots of \(P\), counted with inseparable multiplicity,
have absolute value \(s\), and the elementary symmetric functions give
\begin{equation}\label{eq:p062-annulus-endpoint-coefficient-bounds}
                  |c_i|\leq s^{d-i},
                  \qquad |c_0|=s^d=:S.
\end{equation}
If \(|T|<s\), the constant term strictly dominates every other term of
\(P\); if \(|T|>s\), the leading term strictly dominates.  Therefore
\begin{equation}\label{eq:p062-annulus-endpoint-open-tube}
 X_y:=\{|P(T)|<S\}\subseteq\{|T|=s\}\subseteq A.
\end{equation}
Moreover \(X_y\) is ambient open in the analytic torus.  Thus it is contained
in the topological, hence absolute analytic, interior of \(A\).  Notice that
no assertion that \(X_y\) lies in one specified residue branch is needed.

The polynomial \(P\) induces a finite flat morphism of degree \(d\)
\begin{equation}\label{eq:p062-annulus-endpoint-open-map}
 p\colon X_y\longrightarrow D^-_S:=\{|Z|<S\},
                  \qquad Z=P(T).
\end{equation}
Indeed, it is the base change of the finite free polynomial morphism, with
basis \(1,T,\ldots,T^{d-1}\).  Let
\[
                         H=\operatorname{Tr}_p(h|_{X_y}),
\]
where the fibre trace uses analytic local degrees.  The finite-flat
continuous-trace theorem and
\eqref{eq:p062-annulus-endpoint-zero-current} make \(H\) continuous with
\(
                    \mathrm d'\mathrm d''[H]=0
\)
on the standard strict open disc \(D^-_S\).  The arbitrary-field strict-open
disc theorem makes \(H\) constant.  The zero fibre is the single analytic
point
\(
                \mathcal M(K[T]/(P))=\{y\}
\)
with local degree \(d\), and hence
\begin{equation}\label{eq:p062-annulus-endpoint-trace-at-zero}
                              H(0)=d\,h(y).
\end{equation}

We next identify the limiting fibre, without defining a closed trace of
\(h\).  The closed polynomial tube is
\begin{equation}\label{eq:p062-annulus-endpoint-closed-tube}
 \overline X_y:=\{|P(T)|\leq S\}=\{|T|\leq s\}.
\end{equation}
The equality follows directly from
\eqref{eq:p062-annulus-endpoint-coefficient-bounds} and the same strict
dominance argument.  The polynomial extends to a finite flat degree-\(d\)
map
\[
             \overline p\colon\overline X_y\longrightarrow
             D^+_S:=\{|Z|\leq S\}.
\]
Only this finite map and its compact source will be used; the function \(h\)
is not extended to \(\overline X_y\setminus A\).

Let \(\xi_S\) be the boundary Gauss point of \(D^+_S\).  Then
\begin{equation}\label{eq:p062-annulus-endpoint-single-gauss-fibre}
       \overline p^{-1}(\xi_S)=\{\eta_s\},
       \qquad \deg_{\eta_s}(\overline p)=d.
\end{equation}
To see this, choose \(u\in K^\times\) with \(|u|=s\), and normalize
\[
                Q(U)=u^{-d}P(uU),
                \qquad V=u^{-d}Z.
\]
All coefficients of \(Q\) are power-bounded, while its leading and constant
coefficients are units.  Its reduction \(q(U)\) therefore has degree exactly
\(d\), and at the source and target Gauss points the residue-field map is
\[
  \widetilde K(V)\longrightarrow\widetilde K(U),
                  \qquad V\longmapsto q(U).
\]
The element \(q(U)\) is transcendental and
\(
  [\widetilde K(U):\widetilde K(q(U))]=d
\), including the inseparable degree.  The graded fundamental inequality
thus gives local degree at least \(d\) at \(\eta_s\).  Finite flatness makes
the sum of the local degrees in the fibre equal to \(d\), proving
\eqref{eq:p062-annulus-endpoint-single-gauss-fibre}.

For \(0<q<S\), let \(\xi_q\) be the \(q\)-Gauss point of \(D^-_S\).  We
claim that, as \(q\) increases to \(S\), the entire finite fibre
\(\overline p^{-1}(\xi_q)\) eventually lies in every prescribed
neighbourhood of \(\eta_s\) in \(\overline X_y\).  If \(W\) is such a
neighbourhood, its compact complement
\(
                  C=\overline X_y\setminus W
\)
has compact, hence closed, image under the finite map \(\overline p\).
By \eqref{eq:p062-annulus-endpoint-single-gauss-fibre}, this image does not
contain \(\xi_S\).  Since \(\xi_q\to\xi_S\), it follows that
\(
                  \overline p^{-1}(\xi_q)\subset W
\)
for all sufficiently large \(q<S\).

Every point above such a \(\xi_q\) satisfies \(|P|=q<S\), and hence belongs
to \(X_y\subset A\), where \(h\) is defined.  Given \(\varepsilon>0\), use
continuity of \(h\) at \(\eta_s\) to choose \(W\) so that
\(
        |h-h(\eta_s)|<\varepsilon
\)
on \(W\cap A\).  Since the total calibrated degree of every fibre is \(d\),
the preceding compact-complement convergence gives
\[
             |H(\xi_q)-d\,h(\eta_s)|\leq d\varepsilon
\]
for all sufficiently large \(q<S\).  Therefore
\begin{equation}\label{eq:p062-annulus-endpoint-trace-limit}
                         H(\xi_q)\longrightarrow d\,h(\eta_s).
\end{equation}
This is a direct endpoint limit.  It neither evaluates \(h\) outside \(A\)
nor invokes a closed trace of \(h\).

The function \(H\) is constant on \(D^-_S\).  Combining
\eqref{eq:p062-annulus-endpoint-trace-at-zero} and
\eqref{eq:p062-annulus-endpoint-trace-limit} yields
\begin{equation}\label{eq:p062-annulus-endpoint-rigid-value}
                              h(y)=h(\eta_s).
\end{equation}
Rigid points are dense in a Hausdorff strict analytic space over a
nontrivially valued field
\cite[Proposition~2.1.15]{Berk12}; applying this locally on \(A\), every
endpoint point is a limit of rigid endpoint points.  Hence continuity,
\eqref{eq:p062-annulus-endpoint-core-formula}, and
\eqref{eq:p062-annulus-endpoint-rigid-value} prove
\eqref{eq:p062-annulus-endpoint-affine} on all of \(A\).

Finally, \(r<R\), so \(\log|T|\) is nonconstant on the annular skeleton.
This proves uniqueness of \(a,b\).
\end{proof}

% Quarantined v4 source-correct strict-open dependency for P06.2.
% The legacy fibre-slicing proof has been deleted, rather than patched in
% place.  The active polycylinder proof cites the principal v4 proposition,
% whose local moment step uses CLD Proposition 15.1.1(b) in the forward
% direction.  These new exact bytes are not production-verified.

\begin{proof}[Proof of \cref{prop:p062-arbitrary-field-annulus-v4-source-correct}]
For \(s\in(r,R)\), let \(\eta_s\) be the \(s\)-Gauss point. Testing the
current equation on compactly supported functions of
\(-\log|T|\), and using the weight-one split-torus integration formula,
shows that
\begin{equation}\label{eq:p062-annulus-radial-affine-v4-source-correct}
             u\longmapsto h(\eta_{\mathrm e^{-u}})
\end{equation}
has zero distributional second derivative on
\((-\log R,-\log r)\). Thus it is affine. This is exactly the
one-dimensional part of the split-torus calculation used in
\cref{prop:p062-rational-centred-traces}; unlike the disc case, there is
no endpoint at \(T=0\) and its affine slope is not forced to vanish.

We next prove that every rigid point has the value prescribed by
\eqref{eq:p062-annulus-radial-affine-v4-source-correct}. Let \(y\) be rigid, put
\(E=\mathscr H(y)\), \(s=|T(y)|\), and let
\[
                 P(T)=T^d+\sum_{i=0}^{d-1}c_iT^i
\]
be the monic irreducible polynomial of \(T(y)\). Completeness of \(K\)
gives a unique extension of its absolute value to every finite extension.
Every conjugate of \(T(y)\), counted with inseparable multiplicity, has
absolute value \(s\). Consequently
\begin{equation}\label{eq:p062-annulus-coefficient-bounds-v4-source-correct}
             |c_i|\leq s^{d-i},
             \qquad |c_0|=s^d=:S\in|K^\times|.
\end{equation}
Strict dominance of the constant term for \(|T|<s\) and of the leading
term for \(|T|>s\) gives
\begin{equation}\label{eq:p062-annulus-polynomial-tube-v4-source-correct}
 X_y:=\{|P(T)|<S\}\subseteq\{|T|=s\}\subset A^-_{r,R}.
\end{equation}
The polynomial \(P\) induces a finite flat morphism of degree \(d\)
\[
             p\colon X_y\longrightarrow D^-_S,
             \qquad Z=P(T).
\]
Let \(H=\operatorname{Tr}_p(h|_{X_y})\), with analytic local degrees in
the fibre trace. By \cref{prop:p062-finite-flat-continuous-trace}, \(H\)
is continuous and \(\mathrm d'\mathrm d''[H]=0\). The arbitrary-field
strict-disc theorem \cref{prop:p062-arbitrary-field-disc} makes \(H\)
constant. The zero fibre is the single analytic point \(y\), with degree
\(d\), so
\begin{equation}\label{eq:p062-annulus-trace-zero-v4-source-correct}
                              H=d\,h(y).
\end{equation}

Let \(\xi_S\) denote the boundary Gauss point of the target closed disc.
The graded Gauss calculation used in the proof of
\cref{prop:p062-arbitrary-field-disc}, now at \(t=s\), gives
\begin{equation}\label{eq:p062-annulus-single-gauss-fibre-v4-source-correct}
                      P^{-1}(\xi_S)=\{\eta_s\},
                      \qquad \deg_{\eta_s}(P)=d.
\end{equation}
Indeed, the initial polynomial
\(\operatorname{in}_s(P)(U)\) has degree \(d\); the rational-function
extension generated by \(U\) over that generated by
\(\operatorname{in}_s(P)(U)\) has degree \(d\), including its inseparable
part. The graded fundamental inequality gives local degree at least \(d\),
and finite flatness leaves neither degree nor another point in the fibre.

Choose a neighbourhood \(W\) of \(\eta_s\) whose closure is contained in
the annulus. Since \(P\) is finite and
\eqref{eq:p062-annulus-single-gauss-fibre-v4-source-correct} is a singleton fibre,
there is an analytic neighbourhood \(V\) of \(\xi_S\) such that
\(P^{-1}(V)\subset W\). Applying the continuity part of
\cref{prop:p062-finite-flat-continuous-trace} to
\(P^{-1}(V)\to V\) shows, as \(q\uparrow S\), that
\[
             H(\xi_q)\longrightarrow d\,h(\eta_s).
\]
Together with \eqref{eq:p062-annulus-trace-zero-v4-source-correct}, this yields
\begin{equation}\label{eq:p062-annulus-rigid-to-skeleton-v4-source-correct}
                              h(y)=h(\eta_s).
\end{equation}
Rigid points are dense locally on the strict annulus. The affinity of
\eqref{eq:p062-annulus-radial-affine-v4-source-correct},
\eqref{eq:p062-annulus-rigid-to-skeleton-v4-source-correct}, and continuity give
\eqref{eq:p062-annulus-affine-v4-source-correct} at every analytic point.
Uniqueness is immediate because \(\log|T|\) is nonconstant on the
skeleton.
\end{proof}

\begin{proof}[Proof of \cref{thm:p062-arbitrary-field-polycylinder-v4-source-correct}]
We first work with the open polycylinder. Apply
\cref{prop:p062-strict-product-fibre-slicing-v4-research} to any coordinate
projection. On a disc fibre, \cref{prop:p062-arbitrary-field-disc} shows
that the restriction of \(h\) is constant. Hence \(h\) is independent of
every disc coordinate.
If \(q=0\), this already proves that \(h\) is constant, and there is
nothing more to show.

Fixing the \(K\)-rational origin in every disc factor and replacing the
disc coordinates successively shows, by the preceding constancy, that
\(h\) is the pullback of a continuous function on the product of the
annulus factors. We continue to denote that function by \(h\).

On an annulus fibre, \cref{prop:p062-arbitrary-field-annulus-v4-source-correct}
shows that the restriction of \(h\) is affine in the logarithmic radius.
In particular, it has at an arbitrary fibre point the same value as at the
Gauss point with the same radius. The slicing proposition is valid over
every base point, with its possibly nonrigid complete residue field. It
remains to verify that successive replacements in distinct coordinates
have a common, order-independent endpoint.

Fix radii \(\mathbf s=(s_1,\ldots,s_q)\) in the open radius rectangle. For
a Laurent polynomial
\(
 f=\sum_{\nu\in\mathbf Z^q}a_\nu S^\nu
\), successive one-variable Gauss extensions, in any order, give
\begin{equation}\label{eq:p062-iterated-gauss-v2}
                  |f|_{\mathbf s}
                   =\max_{\nu}|a_\nu|
                         \prod_{j=1}^q s_j^{\nu_j}.
\end{equation}
Indeed, each step takes the maximum over one exponent and the finitely
many maximum operations commute. At an intermediate step the coefficients
may lie in the completed residue field of the coordinates not yet
replaced. The one-variable Gauss norm is still defined by the same maximum
formula, and passing to that completion does not change the result.

The formula extends from Laurent polynomials to every Laurent series
converging on a fixed compact subannulus. Such a series is the Banach-norm
limit of finite Laurent sums, and every bounded multiplicative seminorm is
continuous for that norm. It follows that every order of replacement
induces the same seminorm on each compact subannulus, hence the same
analytic point on the open polyannulus. By
\eqref{eq:p062-iterated-gauss-v2}, that point is the canonical
multivariate Gauss point \(\sigma(\mathbf s)\).

The section \(\mathbf s\mapsto\sigma(\mathbf s)\) is continuous. Indeed,
on a compact logarithmic subrectangle, the norm of a Laurent polynomial is
the maximum of finitely many continuous monomials in \(\mathbf s\), and
the preceding Banach approximation gives the norm of every analytic
function as a locally uniform limit of such functions. This is precisely
continuity for the Berkovich topology.

Now let \(z\) be an arbitrary point of the open polyannulus and put
\(s_j=|S_j(z)|\). Replacing the coordinates one at a time does not change
the value of \(h\), while the preceding calculation identifies the final
point independently of order. Therefore
\[
                     h(z)=h\bigl(\sigma(\mathbf s)\bigr).
\]
It follows that \(h\) factors through the standard tropicalization:
\begin{equation}\label{eq:p062-polycylinder-factor-v4-source-correct}
 h=\phi\circ F,
 \qquad
 F=(-\log|S_1|,\ldots,-\log|S_q|),
\end{equation}
where
\[
             \phi(\mathbf u)
                =h\bigl(\sigma(\mathrm e^{-u_1},\ldots,
                                      \mathrm e^{-u_q})\bigr).
\]
The continuity just proved for the Gauss section and continuity of \(h\)
show that \(\phi\) is continuous on the open rectangle
\(
 \prod_j(-\log R'_j,-\log r_j).
\)

If \(p>0\), choose a compactly supported smooth positive top-degree form
\(\gamma\) in the product of the disc factors with
\(\int\gamma>0\). If \(p=0\), put \(\gamma=1\) and
\(\int\gamma=1\). For an arbitrary compactly supported Lagerberg form
\(\beta\) of bidegree \((q-1,q-1)\) on the logarithmic rectangle, use as
a test form the product of \(\gamma\) with the pullback of \(\beta\).
Since \(\gamma\) has top degree in the disc variables,
\(\mathrm d'\gamma=\mathrm d''\gamma=0\). The weight-one split-torus
integration formula and Fubini therefore convert
\(
 \mathrm d'\mathrm d''[h]=0
\)
into
\(
 \mathrm d'\mathrm d''[\phi]=0
\)
in the distributional Lagerberg sense on that rectangle; the nonzero
factor \(\int\gamma\) cancels. Equivalently,
every distribution
\(
 \partial^2\phi/\partial u_i\partial u_j
\)
vanishes. Convolution on relatively compact Euclidean balls makes every
regularization affine; compatibility on overlaps and continuity show that
\(\phi\) itself is affine. This proves
\eqref{eq:p062-polycylinder-affine-v4-source-correct}.

Each \(\log|S_j|\) is the logarithm of the norm of an analytic unit.
Thus the displayed function is ordinary CLD-smooth with zero Hessian,
weakly smooth and Gubler--Jell--Rabinoff harmonic, and classically
pluriharmonic. Uniqueness follows by restriction to the standard
polyannulus skeleton.
\end{proof}

Proposition 15.1.1. — Soit X un espace k-analytique topologiquement séparé, consi-
déré comme espace G-tropical affiné.
   a) Soit ω une G-forme de type (p, q) sur X. Il existe un G-recouvrement (Vi ) de X et,
pour tout i, un moment fi : Vi → Ti sur Vi , une forme αi de type (p, q) sur une partie
                                                                 ∗
paralinéaire de Ti,trop contenant fi,trop (Vi ) tels que ω|Vi = fi,trop αi .
   b) Soit ω une forme de type (p, q) sur X. Il existe un recouvrement ouvert (Vi ) de X
et, pour tout i, un moment fi : Vi → Ti sur Vi , une forme αi de type (p, q) sur une partie
                                                                 ∗
paralinéaire de Ti,trop contenant fi,trop (Vi ) tels que ω|Vi = fi,trop αi .

# Dependency manifest

The active chain is:

`endpoint theorem -> strict core theorem -> principal v4 fibre slicing -> CLD Proposition 15.1.1(b)`.

The strict annulus branch additionally uses the materialized rational-centred,
finite-trace, and arbitrary-field disc dependencies.  Endpoint annuli use the
separate endpoint packet.  All exact files under `dependencies/` are covered
by `SHA256SUMS`.  The old strict-open fibre-slicing proof is not present.

Production provenance is deliberately separated from mathematical use: copied
receipts/responses authenticate only the exact statement/proof hashes named in
those records.  The new strict-open v4 dependency is quarantined pending an
independent audit and production verification.

# Source matrix

| Use | Exact source | Direction/status |
|---|---|---|
| Fibre local moment presentation | CLD Proposition 15.1.1(b), lines 21482--21491 of `/private/tmp/CLD25.txt` | forward: an ordinary smooth form locally admits a moment presentation |
| Fibre slicing | principal v4 proposition and proof | active; no Corollary 15.1.2 dependency |
| Strict-open annulus | `strict_open_source_correct_statement.tex` and `strict_open_source_correct_proof.tex` | inherited proof block; exact v4 bytes quarantined |
| Strict-open polycylinder | same two files | active proof cites principal v4 slicing |
| Rational-centred/descent | `p062_rational_centred_and_descent_*_verified.tex` | production-verified exact dependency |
| Finite trace/open disc | `p062_finite_trace_and_disc_*_verified.tex` | production-verified exact dependency |
| Endpoint annuli | `endpoint_annulus_*` | production-verified; used only for oc/co/cc |

CLD Corollary 15.1.2 is not used to infer a local moment presentation.

# Independent exact-byte hostile re-audit: P06.2 endpoint polycylinder v4

Date: 30 August 2026.

## Verdict

**PASS as a frozen, quarantined research candidate.**

- Critical errors: **0**.
- Mathematical/proof gaps: **0**.
- Source/dependency gaps: **0**.
- Frozen-byte mismatches: **0**.
- Compile failures: **0**.
- Undefined references/citations: **0**.
- Nonfatal compile diagnostics: **3 underfull boxes in audit wrappers**.
- Visual legibility defects: **0**.

This verdict does not constitute production verification or authorize live-book
integration.  The new source-correct strict-open bytes are explicitly
quarantined pending production review.

## Exact target and integrity

Frozen dossier:

`/private/tmp/nappt_p071_20260827/manuscript/editorial/solver_dossiers_20260830/p062_endpoint_polycylinder_rigid_slicing_v4_20260830/`

| Artifact | SHA-256 | Result |
|---|---|---|
| principal statement | `56ffaad6c42e0a14c56279e263d326120ce03f6ae132d8261deee56ff9ca7663` | match |
| principal proof | `2616ac69cd450010b8f0d108d030bd69d2ff2728cad9641ad7d9d0e1a981a0b7` | match |
| active strict-open statement | `44d43747e47b07c141e1dfbe3932cca29ef24ea0060aee0edfb7529c182f8605` | match |
| active strict-open proof | `9db508b98101b1976af1c14c963e2352aa008a7c440d94e6fe76d7e92e049ecd` | match |
| principal wrapper | `0f234329b64519f3012e5e4e3cfd31911794d9afc73c5d1abbe1de8ce165ae0d` | match |
| principal frozen PDF | `02b840eab83a42435db3428c71ee6a71c5f6e59659af470300c9da1f308a147f` | match |
| strict wrapper | `17bfb9b3cc6ac09e5b1e831457ecdf9149dce67fc9a56d1005ec72db054b6a62` | match |
| strict frozen PDF | `d96e7b1f6b76c0c77e5568a27902faa910df922b2356c2d2b219f1d5890e902a` | match |
| `SHA256SUMS` | `45aeebfbd055a79bb5eca5d20d167df1c91caa3453d375d410b222157f3b1f17` | observed |

Independent `shasum -a 256 -c SHA256SUMS` checks returned `OK` for every
listed artifact before and after the audit.

## Active dependency graph and deletion of the legacy slicing block

**PASS.**

The active mathematical chain is exactly

`endpoint theorem -> strict-core theorem -> principal v4 fibre slicing -> CLD Proposition 15.1.1(b)`.

The active strict-open statement contains only:

1. `prop:p062-arbitrary-field-annulus-v4-source-correct`, lines 6--19; and
2. `thm:p062-arbitrary-field-polycylinder-v4-source-correct`, lines 21--53.

The active strict-open proof contains only the corresponding proof blocks,
lines 7--89 and 91--202.  The former legacy label
`prop:p062-product-fibre-zero-current-candidate`, its statement, and its proof
have no occurrence in either active strict-open file.  There is likewise no
occurrence of `Corollary~15.1.2` in any active TeX dependency.

The strict polycylinder proof, lines 91--202, invokes only the principal v4
slicing proposition at lines 92--96.  Its exact range hash is
`4da2cd94de8dae6210dd616c896fd6277f592e31492abb213b732a81fabaf414`.
The resulting dependency is not circular: the principal slicing proposition
and proof are independent of both strict-open classifications, while only the
later principal endpoint theorem imports them.

## CLD local-moment direction

**PASS.**

The principal slicing proof, lines 37--47, starts with an ordinary smooth form
and invokes CLD Proposition 15.1.1(b) to produce an ordinary-open local moment
presentation.  The range hash is
`46ae9ed416d648a6cfbe1977a521902e106410f3990cd98f6e428d4bb0f57667`.

The locally materialized extraction
`dependencies/CLD25_Proposition_15_1_1_b_lines_21482_21491.txt`, SHA-256
`d2c41ff1bb5b3698e9124458932f20fa20956d70bdcae790320ad8656ddd1f7a`,
is byte-identical to lines 21482--21491 of `/private/tmp/CLD25.txt`, SHA-256
`021be260392de1f3240858d4c775da86653518dfab4fd9c4c21a2406081b6c7e`.
Clause (b) states precisely the direction used.  Corollary 15.1.2 is not used
to infer a moment presentation anywhere in the active chain.

## Fully open and endpoint annulus routing

**PASS.**

The principal endpoint proof, lines 276--284, makes a disjoint four-way split:

- `r < |S| < R` uses only
  `prop:p062-arbitrary-field-annulus-v4-source-correct`;
- open--closed, closed--open, and closed--closed use only
  `prop:p062-annulus-endpoint-arbitrary-field`.

The strict-core theorem at lines 295--302 similarly uses only the
source-correct strict-open polycylinder theorem.  The combined routing range
has SHA-256
`a53a0457fd144ae03bdf238fadcf23e8abaa5eef26bc9b21191478762219d028`.

The endpoint statement, SHA-256
`a54723f040926d5a309494bb77d3c734c1c9d27bfb9fd18c56fc70ea4d03c9e1`,
lists exactly the three endpoint patterns.  Its proof and production response
hashes are
`461afad5804a3269cb385550af6f4abcb5543d4f3151db56f986490c74e6947e`
and
`6b5aa2445bbd8b1dbe594fc23571c979d64eb5dc839f038101ce97b93181cee4`.
The receipt records verdict `correct`, zero critical errors, and zero gaps for
that exact scope.

The active fully open annulus proof is the previously verified mathematical
block with only labels adjusted to the new source-correct namespace.  Its
lines 7--89 have range SHA-256
`5ac61c2e1f7aa7331646f5c7c085a371e69fddc387618bfeca371565fffd77d6`.
It does not import the deleted slicing proposition.

## Rational-centred, descent, finite trace, and disc dependencies

**PASS.**

The materialized rational-centred/descent statement and proof have SHA-256
values

- `7510083fecda2d561c4ce2a1b35d08377268ef48adda2b242172428dbd39caec`;
- `77a013d4e655d288d8803c9650df5446bed570de2644591edddfc116c323eab2`.

They are byte-identical to the frozen release-candidate originals.  The copied
production response has SHA-256
`112168020192b4dffb65d309898992528dac0c8bdf8ba3aeddfe3bb245597c41`
and verdict `correct`, with zero errors, gaps, and reference issues.

The materialized finite-flat trace/disc statement and proof have SHA-256
values

- `67a8d8a21e23cb34cf9b8f6685efbf2c9eaa6ff484b3696cbc1867a8940e095c`;
- `479643ee3aa534bbfc98adcf69837356d3cd8eff22d92accd1319aabb049553e`.

They are byte-identical to the production materialization and curve-candidate
originals.  The dependency-closed response has SHA-256
`1459e3e5b05666c9537491e945b96b159df4448c502a9c4c930bdd9a204d11ea`
and verdict `correct`, with zero critical errors and gaps.

The strict annulus proof uses these dependencies in their verified directions:
rational-centred split-torus calculation, finite-flat continuous trace, and
open-disc constancy.  The polynomial tube, singleton full-degree Gauss fibre,
inseparable multiplicity, and rigid-density passage are the same verified
mathematical steps; no endpoint conclusion is imported into the fully open
proof.

## Fibre slicing, ordinary neighbourhoods, and current restriction

**PASS.**

- Principal proof lines 14--20 choose a strict affinoid neighbourhood with the
  point in its relative interior and then an ordinary open neighbourhood with
  compact closure in that interior.
- Lines 27--35 place all later supports in the ordinary relative interior of
  the product domain.
- Lines 37--89 construct an exact smooth ambient spread of every compactly
  supported fibre test, using the correct CLD moment direction, equality of
  fibre/ambient completed residue fields, density, ultrametric equality, and a
  finite smooth partition.
- Lines 90--95 extend by zero only after compact support is contained in an
  ordinary relative interior.  No arbitrary current is pulled back through an
  affinoid boundary.
- Lines 97--152 insert a compact cutoff before smooth approximation and justify
  both sides of the continuous-multiplier Fubini limit.
- Lines 154--178 pair the original current only with a legitimate global
  compactly supported smooth test form.  The bidegree calculation includes
  base dimension zero.

Thus the principal slicing proposition supplies exactly the restriction used
by the active strict polycylinder proof, for arbitrary (including nonrigid)
base points.

## Endpoint combinations, finite scalar extension, and dimension zero

**PASS.**

- The principal statement allows nonnegative `p,q` and fixes empty product and
  empty sum conventions.
- Proof lines 181--183 dispose of `p=q=0` before choosing a factor.
- Lines 185--217 construct the coordinatewise interior `Q`, prove
  `Q subset Int(P)`, and prove density via rigid points.
- Lines 219--232 explicitly identify an empty complementary product with
  `M(K)`; slicing therefore includes a zero-dimensional base.
- Lines 234--263 use only finite residue-field extensions when identifying the
  scalar extension of a factor interior with the absolute interior of the
  scalar-extended factor.  No arbitrary complete-extension equality is
  claimed.
- Lines 265--284 apply the disc, fully open annulus, and three endpoint
  classifications with the exact routing audited above.
- Lines 304--349 induct on every included endpoint coordinate.  The disc zero
  section and two finite-extension interior annulus radii reduce the induction
  count, including all mixtures of open, closed, and half-closed factors.
- Lines 351--358 use rigid density and continuity and prove uniqueness for both
  `q=0` and `q>0`.

The active strict-open theorem itself also handles empty factors: if `q=0`,
disc-fibre constancy gives the result (and `p=q=0` is the one-point product);
if `p=0`, the test factor is explicitly set to `gamma=1` with integral one.

No counterexample or missing endpoint/dimension-zero case was found.

## V3-to-v4 exact comparison

After replacing version labels by a common placeholder, the principal
statements are byte-identical.  The normalized principal proof differs only
by routing the fully open annulus and strict core from the old dependency
labels to the source-correct v4 labels and by accurately calling those new
bytes quarantined rather than production-verified.

The active strict-open statement deletes the old fibre-slicing proposition
entirely.  The active strict-open proof deletes its old proof entirely and
changes the strict polycylinder proof's sole slicing call to the principal v4
proposition.  No other mathematical scope is enlarged.

## Independent compilation and visual audit

All required sources and dependencies were copied without editing to
`/private/tmp/p062_v4_audit.M32JmM/`.  Both wrappers were compiled there with

`latexmk -pdf -interaction=nonstopmode -halt-on-error`.

Results:

| Wrapper | Pages | Independent PDF SHA-256 | Errors/undefined refs |
|---|---:|---|---:|
| principal | 8 A4 | `2df94308cada0ee3b236a974f7f3ad560dda03ca99e9ea249067f3f554adf1c8` | 0 |
| active strict dependency | 7 A4 | `45f7c91c38f914a799f6a2169df842486e5bac64ea236a59322b1933b0807cec` | 0 |

The final principal log contains one nonfatal `Underfull hbox` in the
dependency-placeholder paragraph.  The final strict log contains two
nonfatal `Underfull hbox` diagnostics in justified dependency statement text.
There are no overfull boxes, undefined references/citations, multiply-defined
labels, fatal errors, or package errors.

All fifteen pages were rendered at 110 dpi and inspected.  The underfull
paragraphs cause only loose word spacing in audit-only wrapper text.  All
theorem statements, proofs, equations, citations, margins, page breaks, and
bibliographies remain legible, with no clipping, collision, missing glyph,
black rectangle, or broken reference.  The three diagnostics are therefore
recorded but are not a mathematical, source, or compile-failure gate.

## Non-actions and scope

No frozen dossier file and no live-book file was edited.  No production or
research verification service was called.  This report is the only persistent
audit artifact created.

The verdict is limited to the stated standard strict factors and the exact
source/dependency chain.  It does not assert a result for non-strict radii,
arbitrary analytic germs, or arbitrary complete scalar extensions in the
endpoint-interior identity.

% Quarantined v4 research proof.  The finite-residue-field step is
% deliberate: no equality of analytic interiors is asserted after an
% arbitrary complete scalar extension.

\begin{proof}[Proof of
\cref{prop:p062-strict-product-fibre-slicing-v4-research}]
Write \(\pi\colon X\to B\) for the projection.  It is flat,
boundaryless, and of pure relative dimension one.  Fix \(x\in B\) and a
compactly supported ordinary CLD-smooth function \(\varphi\) on
\(
 X_x=Z_{\mathscr H(x)}
\).

Since \(B\) is good and strict, choose a strict affinoid neighbourhood
\(U_0\) of \(x\) in \(B\) such that
\(
 x\in\operatorname{Int}(U_0/B)
\).  Since \(B\) is locally compact and Hausdorff, there is an ordinary
open neighbourhood \(V\) of \(x\) whose closure is compact and contained
in \(\operatorname{Int}(U_0/B)\).  Put
\[
                         Y_0=U_0\times_KZ,
                         \qquad F=(Y_0)_x=X_x.
\]
The space \(Y_0\) is good and topologically separated, hence almost good
and locally holomorphically separated.  Moreover,
\begin{equation}\label{eq:p062-endpoint-relative-interior-support-v4-research}
 V\times_KZ\subseteq\operatorname{Int}(Y_0/X),
\end{equation}
because \(V\times_KZ\) is an ordinary open subset of \(X\) contained in
the analytic domain \(Y_0\); here we use the identification of relative
interior with topological interior for an analytic domain
\cite[Proposition~1.5.5(i)]{Berk93}.  We spread \(\varphi\) to a smooth
function on \(Y_0\), exactly on the fibre, and keep its support inside this
ordinary open subset of \(X\).

Fix \(y\in\operatorname{supp}(\varphi)\).  The ordinary-open local moment
presentation of an ordinary form
\cite[Proposition~15.1.1(b)]{CLD25} gives a neighbourhood \(V_y'\) of \(y\) in
\(F\), invertible \(\mathscr H(x)\)-analytic functions
\(f_{y,1},\ldots,f_{y,N_y}\) there, and a smooth real function \(\rho_y\)
such that
\[
 \varphi=\rho_y\bigl(-\log|f_{y,1}|,\ldots,
                         -\log|f_{y,N_y}|\bigr)
 \qquad\text{on }V_y'.
\]
The completed residue field of \(y\), computed in the analytic fibre
\(F\), is canonically the same as that computed in \(Y_0\).  Indeed, on
affinoid charts associated with \(A\to C\), adjoining the image of
\(\mathscr H(x)\) to \(C\) does not change the completion of the fraction
field of the image of \(C\): the latter is already dense in
\(\mathscr H_{Y_0}(y)\), and the image of \(\mathscr H(x)\) lies in that
closure.

Since \(Y_0\) is good at \(y\), its local residue field is dense in
\(\mathscr H_{Y_0}(y)\) by \cite[\S11.1.10]{CLD25}.  For each \(k\), choose
a local ambient meromorphic residue class represented by an analytic germ
\(g_{y,k}\in\mathcal O_{Y_0,y}\) such that
\begin{equation}\label{eq:p062-endpoint-moment-approximation-v4-research}
             |g_{y,k}(y)-f_{y,k}(y)|<|f_{y,k}(y)|.
\end{equation}
After shrinking around \(y\), every \(g_{y,k}\) is an analytic unit and
the ultrametric inequality gives the literal identities
\begin{equation}\label{eq:p062-endpoint-moment-equality-v4-research}
                    |g_{y,k}|_F=|f_{y,k}|.
\end{equation}
We also shrink the ambient neighbourhood so that it is contained in
\(V\times_KZ\).  Thus the ordinary smooth ambient function
\[
 \Phi_y=\rho_y\bigl(-\log|g_{y,1}|,\ldots,
                         -\log|g_{y,N_y}|\bigr)
\]
restricts exactly to \(\varphi\) on a fibre neighbourhood of \(y\), and is
defined on an ambient neighbourhood contained in \(V\times_KZ\).

The support of \(\varphi\) is compact.  Choose finitely many such ambient
neighbourhoods \(W_1,\ldots,W_N\), all contained in \(V\times_KZ\), and use
\cite[Corollary~15.2.9]{CLD25} to choose nonnegative ordinary smooth
functions \(\theta_i\) on \(Y_0\), compactly supported in \(W_i\), whose
sum is one near \(\operatorname{supp}(\varphi)\).  Each
\(\theta_i\Phi_i\) extends smoothly by zero on \(Y_0\).  The sum
\[
                              \Phi=\sum_i\theta_i\Phi_i
\]
has compact support contained in \(V\times_KZ\) and satisfies
\begin{equation}\label{eq:p062-endpoint-exact-spread-v4-research}
                              \Phi|_F=\varphi.
\end{equation}
By \eqref{eq:p062-endpoint-relative-interior-support-v4-research}, this
compact support is contained in the ordinary interior of \(Y_0\) relative
to \(X\).  We may therefore regard \(\Phi\) as a compactly supported smooth
function on the ordinary open subspace \(V\times_KZ\); extension by zero
to \(X\), as in \cite[Proposition~6.4.5(a)]{CLD25}, is ordinary CLD-smooth.
We use the same symbol for this extension.

Set
\[
                         \tau=\mathrm d'\mathrm d''\Phi.
\]
The total space \(X\) is good and topologically separated, hence locally
holomorphically separated.  Decompose the compactly supported symmetric
form \(\tau\) as
\(
 \tau=\tau^+-\tau^-
\)
with \(\tau^\pm\) positive and compactly supported, using
\cite[Proposition~6.5.2(b)]{CLD25}.  Let
\(
 C=\operatorname{supp}(\tau^+)\cup\operatorname{supp}(\tau^-)
\), a compact subset of \(X\).  By
\cite[Corollary~15.2.5]{CLD25}, choose an ordinary smooth compactly
supported function \(\chi\) on \(X\) that equals one on a neighbourhood of
\(C\).  The function \(\chi h\) is continuous and compactly supported.
Applying \cite[Theorem~15.2.6]{CLD25} with errors tending to zero gives
ordinary smooth compactly supported functions \(h_\ell\) such that
\begin{equation}\label{eq:p062-endpoint-cutoff-approximation-v4-research}
 \lVert h_\ell-\chi h\rVert_{C^0(X)}\longrightarrow0.
\end{equation}
In particular, \(h_\ell\to h\) uniformly on \(C\).

For \(t\in B\), set
\[
                 G(t)=\int_{X_t}h_t\,\tau|_{X_t}.
\]
Fibre integration makes
\(
 t\mapsto\int_{X_t}\tau^\pm
\)
continuous by \cite[Theorem~16.3.3(d)]{CLD25}.  These functions vanish
away from the compact projections of the supports of \(\tau^\pm\), and are
therefore bounded on \(B\).  Consequently
\begin{equation}\label{eq:p062-endpoint-uniform-fibre-limit-v4-research}
 \sup_{t\in B}\left|
   \int_{X_t}(h_\ell-h)\tau\right|
 \leq \lVert h_\ell-h\rVert_{C^0(C)}
       \sup_{t\in B}\left(\int_{X_t}\tau^++\int_{X_t}\tau^-\right)
 \longrightarrow0.
\end{equation}
The smooth fibre integrals are continuous, so \(G\) is continuous.  The
same limit, combined with the projection formula
\cite[Theorem~16.5.4 and Corollary~16.5.5]{CLD25}, gives, for every
compactly supported smooth top-degree form \(\alpha\) on \(B\),
\begin{equation}\label{eq:p062-endpoint-continuous-fubini-v4-research}
 \int_X h\,\tau\wedge\pi^*\alpha
                         =\int_B G\,\alpha.
\end{equation}
Indeed, the projection formula applies to each smooth form
\(h_\ell\tau\); the fibre side converges uniformly by
\eqref{eq:p062-endpoint-uniform-fibre-limit-v4-research}, and the total
space side converges against the fixed compactly supported signed
top-degree measure \(\tau\wedge\pi^*\alpha\).

Since \(\alpha\) has top bidegree on the \(m\)-dimensional base,
\(
 \mathrm d'\alpha=\mathrm d''\alpha=0
\), including when \(m=0\).  Therefore
\[
 \mathrm d'\mathrm d''(\Phi\pi^*\alpha)
       =(\mathrm d'\mathrm d''\Phi)\wedge\pi^*\alpha.
\]
The function \(\Phi\) is now a global compactly supported smooth function,
so \(\Phi\pi^*\alpha\) is a legitimate compactly supported test form on
\(X\).  Pairing
\eqref{eq:p062-general-product-zero-current-v4-research} with this test form
and using \eqref{eq:p062-endpoint-continuous-fubini-v4-research} gives
\(
 \int_B G\alpha=0
\)
for every \(\alpha\).  The strict pure-dimensional space \(B\) is very
rich by \cite[Theorem~16.4.4(b)]{CLD25}; hence continuous functions inject
into \((0,0)\)-currents by \cite[Corollary~7.3.12]{CLD25}.  Thus \(G=0\).
At \(x\), \eqref{eq:p062-endpoint-exact-spread-v4-research} now yields
\[
 0=G(x)=\int_{X_x}h_x\,\mathrm d'\mathrm d''\varphi
       =\bigl\langle\mathrm d'\mathrm d''[h_x],\varphi\bigr\rangle.
\]
This proves \eqref{eq:p062-general-fibre-zero-current-v4-research}.
\end{proof}

\begin{proof}[Proof of \cref{thm:p062-endpoint-polycylinder-v4-research}]
If \(p=q=0\), then \(P=\mathcal M(K)\), and the assertion holds with
\(c=h(*)\).  We henceforth assume \(p+q>0\).

We separate the absolute-interior issue from the affine interpolation.  For
every factor \(F_\nu\) of \(P\), write
\(
 F_\nu^\natural=\operatorname{Int}(F_\nu)
\)
for its absolute analytic interior.  Thus an open factor is unchanged; a
closed disc loses only its endpoint Gauss point; and an annulus loses the
Gauss point at each included endpoint but retains every nonskeletal
endpoint residue branch.  The latter one-dimensional description is the
topological input in the verified endpoint theorem
\cref{prop:p062-annulus-endpoint-arbitrary-field}.

Let \(M\) be the product of the ambient affine lines and analytic tori, and
put
\begin{equation}\label{eq:p062-endpoint-coordinate-interior-v4-research}
 Q=\bigcap_\nu\operatorname{pr}_\nu^{-1}(F_\nu^\natural)\subset P.
\end{equation}
Every \(F_\nu^\natural\) is an ordinary open subset of its ambient curve.
Hence \(Q\) is an ordinary open subset of \(M\) contained in \(P\).  It is
therefore contained in the topological interior of the analytic domain
\(P\subset M\).  Berkovich's relative-interior theorem and composition
formula \cite[Proposition~1.5.5(i)--(ii)]{Berk93}, together with the
boundarylessness of \(M\), give
\begin{equation}\label{eq:p062-endpoint-Q-in-interior-v4-research}
                              Q\subset\operatorname{Int}(P).
\end{equation}
Every rigid point of \(P\) lies in \(Q\): its coordinate images are rigid,
whereas the excluded one-dimensional endpoint Gauss points are nonrigid.
Rigid points are dense locally in a strict analytic space
\cite[Proposition~2.1.15]{Berk12}.  Consequently
\begin{equation}\label{eq:p062-endpoint-Q-dense-v4-research}
                              \overline Q=P.
\end{equation}

We next record why slicing at a rigid base point sees the full
one-dimensional absolute interior after passage to the residue field.  Fix
a factor \(F_\nu\), let
\[
 B_\nu=\prod_{\mu\ne\nu}F_\mu^\natural,
 \qquad Z_\nu=F_\nu^\natural.
\]
Here \(B_\nu=\mathcal M(K)\) when the displayed product is empty.  The
product \(B_\nu\times_KZ_\nu\) is \(Q\), so the current equation holds
there by \eqref{eq:p062-endpoint-Q-in-interior-v4-research}.  The spaces
\(B_\nu\) and \(Z_\nu\) are good, topologically separated, pure and strict,
and \(Z_\nu\) is boundaryless.  Thus
\cref{prop:p062-strict-product-fibre-slicing-v4-research} applies, including
when \(B_\nu\) has dimension zero.

Let \(x\in B_\nu\) be rigid and set \(E=\mathscr H(x)\).  Then \(E/K\) is
finite.  For the scalar-extension map
\(
 p_E\colon F_{\nu,E}\to F_\nu
\), one has
\begin{equation}\label{eq:p062-endpoint-finite-interior-equality-v4-research}
       (F_\nu^\natural)_E
       =\operatorname{Int}(F_{\nu,E}/E).
\end{equation}
Indeed, \(p_E\) is finite and hence boundaryless.  Apply the equality case
of Berkovich's composition formula first to
\[
 F_{\nu,E}\longrightarrow F_\nu\longrightarrow\mathcal M(K)
\]
and then to
\[
 F_{\nu,E}\longrightarrow\mathcal M(E)\longrightarrow\mathcal M(K).
\]
The two finite arrows have empty relative boundary, and the remaining maps
are good and locally separated.  The first composition gives
\(
 \operatorname{Int}(F_{\nu,E}/K)
   =p_E^{-1}(\operatorname{Int}(F_\nu/K))
\), while the second identifies its left side with
\(
 \operatorname{Int}(F_{\nu,E}/E)
\).  This is
\eqref{eq:p062-endpoint-finite-interior-equality-v4-research}.  Notice that
this argument uses the finiteness of \(E/K\); no analogous equality is
claimed for an arbitrary complete scalar extension.

Let \(h_x\) be the restriction of \(h\) to the full fibre
\(F_{\nu,E}\).  The slicing proposition and
\eqref{eq:p062-endpoint-finite-interior-equality-v4-research} show that its
Hessian current vanishes on the absolute analytic interior of that fibre.
The verified arbitrary-field one-dimensional results therefore give
\begin{equation}\label{eq:p062-endpoint-one-fibre-classification-v4-research}
 \begin{cases}
 h_x\text{ is constant},&F_\nu\text{ is a disc},\\
 h_x=a_x\log|S_\nu|+b_x,&F_\nu\text{ is an annulus},
 \end{cases}
\end{equation}
where the disc case is
\cref{prop:p062-arbitrary-field-disc}.  For the annulus case, the fully open
domain \(r_\nu<|S_\nu|<R'_\nu\) is covered by the source-correct, quarantined
strict-open result
\cref{prop:p062-arbitrary-field-annulus-v4-source-correct}; the open--closed,
closed--open, and closed--closed domains are covered by
\cref{prop:p062-annulus-endpoint-arbitrary-field}.  Each proposition is
applied over \(E\), and no endpoint proposition is invoked for the fully open
domain.

For each factor let \(F_\nu^-\) denote its coordinatewise strict core:
\[
\begin{aligned}
 F_\nu^-&=\{|T_\nu|<R_\nu\}
              &&\text{for a disc},\\
 F_\nu^-&=\{r_\nu<|S_\nu|<R'_\nu\}
              &&\text{for an annulus},
\end{aligned}
\]
and put \(P^-=\prod_\nu F_\nu^-\).  The source-correct, quarantined strict-open
polycylinder theorem
\cref{thm:p062-arbitrary-field-polycylinder-v4-source-correct} gives constants
\(c,\lambda_1,\ldots,\lambda_q\) for which
\begin{equation}\label{eq:p062-endpoint-core-formula-v4-research}
             h=L:=c+\sum_{j=1}^q\lambda_j\log|S_j|
             \qquad\text{on }P^-.
\end{equation}

We prove \(h=L\) at every rigid point \(y\in P\) by induction on
\[
 e(y)=\#\{\nu:\operatorname{pr}_\nu(y)\notin F_\nu^-\}.
\]
The coordinate image of a rigid point is rigid.  Hence an index counted by
\(e(y)\) is an included endpoint residue branch, never an endpoint Gauss
point.  The case \(e(y)=0\) is
\eqref{eq:p062-endpoint-core-formula-v4-research}.

Suppose \(e(y)>0\), choose a counted coordinate \(\nu\), and let \(x\) be
the projection of \(y\) to the other factors.  Then \(x\in B_\nu\) is
rigid, so \eqref{eq:p062-endpoint-one-fibre-classification-v4-research}
applies over \(E=\mathscr H(x)\).

If \(F_\nu\) is a disc, replace the \(\nu\)-th coordinate by the zero
section of the fibre.  The resulting point \(y_0\) is rigid,
\(e(y_0)=e(y)-1\), and the first line of
\eqref{eq:p062-endpoint-one-fibre-classification-v4-research} gives
\[
                              h(y)=h(y_0).
\]
The function \(L\) has no disc-coordinate term, so the induction hypothesis
gives
\(
 h(y)=h(y_0)=L(y_0)=L(y)
\).

Suppose instead that \(F_\nu\) is an annulus.  Since the valuation on \(E\)
is nontrivial, the divisible hull of \(|E^\times|\) is dense in
\(\mathbf R_{>0}\).  Choose distinct
\[
       \rho_1,\rho_2\in(r_\nu,R'_\nu)
          \cap |E^{\mathrm{alg},\times}|.
\]
For each \(k\), choose a finite extension \(E_k/E\) and an element of
absolute value \(\rho_k\).  It defines a rigid point \(z_k\) of the annulus
fibre, hence a rigid point \(y_k\in P\) above \(x\).  Both points have
\(e(y_k)=e(y)-1\).  By induction,
\begin{equation}\label{eq:p062-endpoint-two-fibre-values-v4-research}
                         h(y_k)=L(y_k)\qquad(k=1,2).
\end{equation}
On this fibre, both \(h_x\) and \(L|_{F_{\nu,E}}\) are affine functions of
\(\log|S_\nu|\).  The two radii are distinct, so
\eqref{eq:p062-endpoint-two-fibre-values-v4-research} makes those affine
functions identical.  In particular \(h(y)=L(y)\).  This completes the
induction.

Rigid points are dense in \(P\), and both \(h\) and \(L\) are continuous.
Thus \eqref{eq:p062-endpoint-polycylinder-affine-v4-research} holds on all
of \(P\).  If \(q=0\), uniqueness of \(c\) follows by evaluating at any
point of the nonempty space \(P\).  If \(q>0\), restrict to the zero
sections in the disc coordinates and to the standard skeleton of the
strict open polyannulus core.  The functions
\(1,\log|S_1|,\ldots,\log|S_q|\) are affinely independent on its open
logarithmic box, which proves uniqueness of all coefficients.

Finally, every \(\log|S_j|\) is the logarithm of the norm of an analytic
unit.  The displayed formula is therefore ordinary CLD-smooth with zero
smooth Hessian and classically pluriharmonic; the already verified
comparison results give weak smoothness and Gubler--Jell--Rabinoff
harmonicity.
\end{proof}
